\documentclass[Total.tex]{subfiles}
\begin{document}
	\chapter{Differentials in Complex Field}
		\section{Previews}
			Review:
			\begin{itemize}
				\item Complex Field
				\item Continuity
			\end{itemize}
			\subsection{Real and Complex Vector Spaces, Complexification}
				Consider the notation as follwing: 
				\begin{itemize}
					\item $Hom_\RR(\CC)$ for $\RR$-linear mappings;
					\item $Hom_\CC(\CC)$ for $\CC$-linear mappings;
				\end{itemize}
				Then,
				\begin{lemma}
					Let $T\in Hom_\RR(\CC)$ be an $\RR$-linear endomorphisms of $\CC$, with following matrix representation for basis s$1,i$
					\begin{gather*}
						T=\begin{pmatrix}
							a&b\\
							c&d
						\end{pmatrix}
					\end{gather*}
					Then, 
					\begin{gather*}
						Tz=\lambda z+\mu \bar{z}\qquad\begin{cases}	\lambda=\lambda_1+i\lambda_2=\frac{1}{2}[(a+d)-(b-c)i]\\
							\mu=\mu_1+i\mu_2=\frac{1}{2}[(a-d)+(b+c)i]
						\end{cases}
					\end{gather*}
					we have: $T$ is $\CC$-linear$\Leftrightarrow a=d,b=-c$.
				\end{lemma}
				\begin{proof}
					With $z=\begin{pmatrix}
						x\\
						y
					\end{pmatrix}$ we have
					\begin{gather*}
						Tz=\begin{pmatrix}
							a&b\\
							c&d
						\end{pmatrix}\begin{pmatrix}
							x\\
							y
						\end{pmatrix}=\begin{pmatrix}
							ax+by\\
							cx+dy
						\end{pmatrix}=(\lambda_1+i\lambda_2)(x+iy)+(\mu_1+i\mu_2)(x-iy)=\begin{pmatrix}
							\lambda_1 x-\lambda_2y+\mu_1x+\mu_2y\\
							\lambda_2 x+\lambda_1 y-\mu_1 y+\mu_2 x
						\end{pmatrix}
					\end{gather*}
					Which implies
					\begin{gather*}
						\begin{cases}
							a=\lambda_1+\mu_1\\
							b=-\lambda_2+\mu_2\\
							c=\lambda_2+\mu_2\\
							d=\lambda_1-\mu_1
						\end{cases}\Rightarrow \begin{cases}
							\lambda_1=\frac{a+d}{2}\\
							\lambda_2=-\frac{b-c}{2}\\
							\mu_1=\frac{a-d}{2}\\
							\mu_2=\frac{b+c}{2}
						\end{cases}\Rightarrow \begin{cases}
							\lambda=\lambda_1+i\lambda_2=\frac{1}{2}[(a+d)-(b-c)i]\\
							\mu=\mu_1+i\mu_2=\frac{1}{2}[(a-d)+(b+c)i]
						\end{cases}
					\end{gather*}
					And $\CC$-Linearity: Notice $id_\CC$ is linear, but $\bar{(-)}:\CC\rightarrow \CC$ is not.
				\end{proof}
		\section{Complex Differentials, Cauchy-Riemann Equation}
			\subsection{Differentiability}
				\subsubsection{Complex Derivative}
					\begin{Def}
						Define
						\begin{gather*}
							f'(z_0):=\lim_{z\rightarrow z_0} \frac{f(z)-f(z_0)}{z-z_0}
						\end{gather*}
					\end{Def}
					Then, we have: $f(z)=u(z)+i(z)$, and notation $u(z)=u(x+iy),v(z)=v(x+iy)$, if $f'(z_0)$ exists,
					\begin{gather*}
						f'(z_0)=\lim_{z\rightarrow z_0} \frac{u(z)-u(z_0)}{z-z_0}+i\lim_{z\rightarrow z_0} \frac{v(z)-v(z_0)}{z-z_0}\\
						\qquad=\lim_{\Delta y\rightarrow 0}\lim_{\Delta x\rightarrow 0} \frac{u(z+\Delta x+i\Delta y)-u(z+\Delta x)}{\Delta x+i\Delta y}\\+\lim_{\Delta x\rightarrow 0}\lim_{\Delta y\rightarrow 0}\frac{u(z+\Delta x)-u(z)}{\Delta x+i\Delta y}\\+i\lim_{\Delta y\rightarrow 0}\lim_{\Delta x\rightarrow 0} \frac{v(z+\Delta x+i\Delta y)-v(z+\Delta x)}{\Delta x+i\Delta y}\\+i\lim_{\Delta x\rightarrow 0}\lim_{\Delta y\rightarrow 0}\frac{v(z+\Delta x)-v(z)}{\Delta x+i\Delta y}\\
						=[-i\frac{\partial u}{\partial y}+\frac{\partial u}{\partial x}-\frac{\partial v}{\partial y}+i\frac{\partial v}{\partial x}](z_0)
					\end{gather*}
					As matrix coordinate
					\begin{gather*}
						\begin{pmatrix}
							\frac{\partial u}{\partial x}(z_0)&-\frac{\partial v}{\partial y}(z_0)\\
							-\frac{\partial u}{\partial y}(z_0)&\frac{\partial v}{\partial y}(z_0)
						\end{pmatrix}
					\end{gather*}
			\subsection{Cauchy-Riemann Equation}
				\begin{theorem}
					(Cauchy-Riemann) 
					\begin{gather*}
						f'(z)\mbox{ is }\CC\mbox{-linear}\Leftrightarrow \begin{cases}
							\frac{\partial u}{\partial x}=\frac{\partial v}{\partial y}\\
							\frac{\partial u}{\partial y}=-\frac{\partial v}{\partial x}
						\end{cases}
					\end{gather*}
				\end{theorem}
				\begin{proof}
					Omitted.
				\end{proof}
			\subsection{Wirtinger Derivation}
				Define:
				\begin{Def}
					With
					\begin{gather*}
						f(z)=f(z_0)+\lambda(z-z_0)+\overline{z-z_0}+o(z-z_0)
					\end{gather*}
					We have
					\begin{gather*}
						\frac{\partial f}{\partial z}(z_0)=\lambda=\frac{1}{2}(\frac{\partial}{\partial x}-i\frac{\partial f}{\partial y})\\
						\frac{\partial f}{\partial \bar{z}}(z_0)=\mu=\frac{1}{2}(\frac{\partial}{\partial x}+i\frac{\partial f}{\partial y})
					\end{gather*}
				\end{Def}
				\begin{proposition}
					We have:
					\begin{itemize}
						\item If $f$　is at $z_0$ (real)-differentiable, then
						\begin{gather*}
							f(z)=f(z_0)+\frac{\partial f}{\partial z}(z_0)(z-z_0)+\frac{\partial f}{\partial \bar{z}}\overline{(z-z_0)}+o(z-z_0)\qquad z\rightarrow z_0
						\end{gather*}
						\item $\frac{\partial}{\partial z}(f\cdot g)=\frac{\partial f}{\partial z}\cdot g+f\cdot \frac{\partial f}{\partial z}$;
						\item We have:
						\begin{gather*}
							\frac{\partial f}{\partial z}=\overline{\frac{\partial \bar{f}}{\partial \bar{z}}}\qquad \frac{\partial \bar{f}}{\partial \bar{z}}=\overline{\frac{\partial \bar{f}}{\partial z}}
						\end{gather*}
						\item For $f:z\mapsto z,g:z\mapsto \bar{z}$ we have
						\begin{gather*}
							\frac{\partial}{\partial z} f=1;\frac{\partial}{\partial \bar{z}}f=0;\frac{\partial}{\partial \bar{z}}g=1;\frac{\partial}{\partial z}f=0
						\end{gather*}
						\item (Chain-Rule I) Suppose $g(w),f(z)$ continuous differentiable, we have
						\begin{gather*}
							\frac{g\circ f}{\partial z}=\frac{\partial g}{\partial w}\circ
						\end{gather*}
						\item (Chain-Rule II) Suppose $\varphi:[a,b]\rightarrow \CC$ continuous differentiable, then
						\begin{gather*}
							\frac{d}{dt}(f\circ \varphi)=\frac{df}{dz}\circ \varphi\cdot \frac{d\varphi}{dt}+\frac{d}{d\bar{z}}\circ \varphi\cdot \overline{\frac{d\varphi}{dt}}
						\end{gather*}
					\end{itemize}
				\end{proposition}
				\begin{proof}
					\begin{itemize}
						\item 
						\item 
						\item 
						\item 
						\item 
						\item 
					\end{itemize}
				\end{proof}
			\subsection{Complex Differentiablity}
				\begin{proposition}
					Suppose $f:U\rightarrow \CC$ be a function with $U\subseteq\CC$ open. The followings are equivalent:
					\begin{itemize}
						\item 
						\item 
						\item 
						\item 
					\end{itemize}
				\end{proposition}
				\begin{proof}
					\begin{itemize}
						\item 
					\end{itemize}
				\end{proof}
				
				\begin{Def}
					(Holomorphic Functions)
					\begin{itemize}
						\item $f$ is \textbf{holomorphic} on $U$, if $f$ is \underline{complex-differentiable} at each $z\in U$;
						\item Take the notion of $\mathcal{O}(U)$ for the collection of holomorphic functions.
					\end{itemize}
				\end{Def}
				
				\begin{proposition}
					(Rules of Complex Derivative) \begin{itemize}
						\item Suppose $f,g$ complex differentiable at $z_0$, then at the point $z_0$, we have
						\begin{gather*}
							内容...
						\end{gather*}
						\item 
					\end{itemize}
				\end{proposition}
				\begin{proof}
					内容...
				\end{proof}
			\subsection{Polynomials,Power Series}
			\subsubsection{Polynomials}
			\begin{proposition}
				Polynomial $p\in \CC[z]$ is complex differentiable, and
				\begin{gather*}
					p'(z)=\sum _{j=1}^n j a_jz^{j-1}
				\end{gather*}
			\end{proposition}
			\begin{proof}
				Omitted.
			\end{proof}
			\subsubsection{Power Series}
			\begin{Def}
				Define the power series:
				\begin{gather*}
					f(z)=\sum_{i=0}^\infty a_i(z-z_0);z_0\in \CC
				\end{gather*}
			\end{Def}
			We have
			\begin{proposition}
				(Cauchy-Hadamard) The convergence radius
				\begin{gather*}
					R=\frac{1}{\lim\sup|c_n|^\frac{1}{n}}
				\end{gather*}
				
				where:
				\begin{itemize}
					\item $|z|>R$ Disvergent;
					\item $|z|=R$ Convergent or Disvergent;
					\item $|z|<R$ Convergent.
				\end{itemize}
			\end{proposition}
			\begin{proof}
				\begin{itemize}
					\item $|z|>R$, 
					\begin{gather*}
						\sqrt[n]{|c_n|}
					\end{gather*}
					\item $|z|<R$,
					\begin{gather*}
						\sqrt{n}{|c_n|}\geq t+\epsilon
					\end{gather*}
					
					for finite of $n$. And for others,
					\begin{gather*}
						\sqrt{n}{c_n}<t+\epsilon\Rightarrow 
					\end{gather*}
				\end{itemize}
			\end{proof}
			\begin{corollary}
				If the radius $R$ and $\lim |\frac{a_{n+1}}{a_n}|$ exists, then,
				\begin{gather*}
					R=\lim |\frac{a_{n+1}}{a_n}|
				\end{gather*}
			\end{corollary}
			Some of review in analysis
			\begin{proposition}
				If $\{a_n\},\{b_n\}$ is absolute convergent, then definen $c_n:=\sum_{k=0}^n a_k b_{n-k}$,
				\begin{gather*}
					\sum_{n=0}^\infty c_n=(\sum a_n)(\sum b_n)
				\end{gather*}
			\end{proposition}
			\begin{proof}
				内容...
			\end{proof}
			
		\section{Fundamental Theorems in Complex Analysis}
			\subsection{Invertible Function}
				\begin{proposition}
					Suppose $f:U\rightarrow V$ continuous \underline{continuous complex differentiable,bijective}. If for every $z\in U$,
					\begin{gather*}
						f'(z)\not=0
					\end{gather*}
					Then, there exists an inverse $f^{-1}:V\rightarrow U$ with 
					\begin{gather*}
						(Df^{-1}(f(z)))=(Df(z))^{-1},z\in U
					\end{gather*}
				\end{proposition}
				\begin{proof}
					内容...
				\end{proof}
				
				\begin{proposition}
					Suppose $f:U\rightarrow V$ be continuous complex differentiable, with $U,V\subseteq \CC$ open. And $f$ is \underline{bijective}, $f'(z)\not=0,z\in U$.
					
					Then, the inverse $f^{-1}:V\rightarrow U$ is also continuous complex differentiable, with
					\begin{gather*}
						(f^{-1})'(w)f'(f(w))=1
					\end{gather*}
				\end{proposition}
				\begin{proof}
					Omitted.
				\end{proof}
			\subsection{Analytic Branch}
				\subsubsection{Logarithms}
					
					\begin{Def}
						Define:
						\begin{itemize}
							\item $C_\alpha=\CC-\{re^{i\alpha}|r\geq 0\},0$
							\item $\Omega_\alpha=\{z\in \CC|\alpha-2\pi<Im z<\alpha\}$
						\end{itemize}
						We have:
						\begin{itemize}
							\item Exponential
							\begin{gather*}
								exp:\Omega_\alpha\rightarrow C_\alpha
							\end{gather*}
							\item Logarithms
							\begin{gather*}
								log_\alpha:C_\alpha\rightarrow \Omega_\alpha
							\end{gather*}
							Especially, logarithm $log_\pi$ is called \textbf{main branch}.
						\end{itemize}
					\end{Def}
				\subsubsection{Roots}
					\begin{Def}
						Define:
						\begin{itemize}
							\item \begin{gather*}
								z\mapsto z^n\qquad \{z\in \CC^*|0\leq |arg(z)|\leq \pi/n\}\rightarrow \CC_\pi
							\end{gather*}
							which is bijective.
							\item  Then, the inverse
							\begin{gather*}
								\sqrt[n]{-}:z\mapsto \sqrt[n]{z}
							\end{gather*}
						\end{itemize}
					\end{Def}
			\subsection{Complex Integration}
				\subsubsection{Area Integraion}
					As what we did in $\RR^2$;
				\subsubsection{Path Integration}
					Define the path integration as following:
					\begin{gather*}
						\int_\gamma f(z)dz=\lim_{\Delta\rightarrow 0} \sum f(z_i)\Delta z_i
					\end{gather*}
					And we have the coordinate decomposition
					\begin{gather*}
						f(z)=u(z)+iv(z)\qquad dz=dx+idy
					\end{gather*}
					Finally:
					\begin{gather*}
						=\int_\gamma (u+iv)(dx+idy)=\int_\gamma udx-vdy+i\int_\gamma vdx+udy
					\end{gather*}
					
					Especially:If the path $\gamma$ is continuous, and $\gamma|[t_i,t_{i+1}]\in C^1([t_i,t_{i+1}])$, we have
					\begin{gather*}
						\int_\gamma f(z)dz=\int_a^b (f\circ\gamma)(t)\cdot \gamma'(t)dt=\sum \int_{t_{i-1}}^{t_i} f(\gamma(t))\cdot \gamma'(t)dt
					\end{gather*}
					\begin{example}
						
					\end{example}
					
					\begin{proposition}
						内容...
					\end{proposition}
					\begin{proof}
						内容...
					\end{proof}
					
				\subsubsection{'Holzhammer Ungleichung'}
					\begin{proposition}
						内容...
					\end{proposition}
					\begin{proof}
						内容...
					\end{proof}
					
					\begin{proposition}
						(Main Theorem of Calculus for Path Integration in Complex Analysis) 
						\begin{gather*}
							\int_\gamma F'(z)dz=F(\gamma(b))-F(\gamma(a))
						\end{gather*}
					\end{proposition}
					\begin{proof}
						内容...
					\end{proof}
			\subsection{The Winding Numbers}
				\begin{Def}
					For a path $\gamma$ with $a\notin sp\gamma$, define
					\begin{gather*}
						wind(\gamma,a)=\frac{1}{2\pi i} \frac{d\gamma}{z-a}
					\end{gather*}
				\end{Def}
				\begin{lemma}
					$f:U\rightarrow \CC$ continuous complex differentiable, with $\gamma$ a path in $U$, then,
					\begin{gather*}
						exp(\int_\gamma \frac{f'}{f}(z)dz)=\frac{f(\gamma(b))}{f(\gamma(a))}
					\end{gather*}
				\end{lemma}
				\begin{proof}
					Consider 
					\begin{gather*}
						\Phi(t)=exp(-\int_a^t \frac{f'}{f}(\gamma(s))\gamma'(s)ds)\cdot f(\gamma(t)) 
					\end{gather*}
					We have:
					\begin{gather*}
						...\\
						\Rightarrow \Phi(a)=\Phi(b),exp(\int_\gamma \frac{f'}{f}(z)dz)=\frac{f(\gamma(b))}{f(\gamma(a))}
					\end{gather*}
				\end{proof}
				\begin{proposition}
					\begin{gather*}
						\oint_{|\gamma-a|=r} \frac{d\xi}{(\xi-z)^n}=\begin{cases}
							0&n\not=-1\\
							2\pi i&n=-1
						\end{cases}
					\end{gather*}
				\end{proposition}
				\begin{proof}
					Now consider $exp()$
				\end{proof}
			\subsection{Cauchy's Formula}
				\subsubsection{Triangle}
					\begin{theorem}
						(Cauchy) If $f$ is analytic on $\Omega\subseteq \CC$, with triangle $\Delta ABC\subseteq \Omega$,
					\end{theorem}
					\begin{proof}
						内容...
					\end{proof}
				\subsubsection{Star Region}
					\begin{Def}
						(Star Region) 
					\end{Def}
				\subsubsection{Stroke$\Rightarrow$ Cauchy's Integration Formula}
			\subsection{Mean Value Property}
				\begin{proposition}
					Suppose $f$ is analytic in $B(z_0,R);R>r$ then,
					\begin{gather*}
						f(z_0)=\frac{1}{2\pi} \int_0^{2\pi} f(z_0+re^{i\theta})d\theta
					\end{gather*}
				\end{proposition}
				\begin{proof}
					内容...
				\end{proof}
				
				\begin{corollary}
					We have
					\begin{gather*}
						|f(z_0)|=\frac{1}{2\pi} \int_0^{2\pi}  |f(z_0+re^{i\theta})|d\theta
					\end{gather*}
				\end{corollary}
				\begin{proof}
					内容...
				\end{proof}
			\subsection{Argument Principle}
			
				Suppsoe that $f(z)=(z-a)^mg(z)$, with $m$ th zero pole, then
				\begin{gather*}
					f'(z)= (z-a)^mg'(z)+m(z-a)^{m-1}g(z);m\geq 1\Rightarrow \frac{f'(z)}{f(z)}=\frac{g'(z)}{g(z)}+\frac{m}{z-a} 
				\end{gather*}
				
				Conversely, suppose $f(z)=(z-a)^{-m}g(z)$ with pole of order $m$, then
				\begin{gather*}
					f'(z)=\Rightarrow 
				\end{gather*}
				
				\begin{theorem}
					(Argument Principle) Let 
					\begin{gather*}
						\frac{1}{2\pi}\int\frac{f'(z)}{f(z)}dz=\sum n(\gamma;z_i)-\sum n(\gamma,p_i)
					\end{gather*}
				\end{theorem}
				\begin{proof}
					内容...
				\end{proof}
			\subsection{Rouche's Theorem}
				\begin{theorem}
					(Rouche) Suppose $D$ be a region on $\CC$, with a 
				\end{theorem}
			\subsection{Open Mapping's Theorem}
				\begin{theorem}
					Suppose $f$ is a non-constant analytic function on the region $\Omega$, then its image
					\begin{gather*}
						f(\Omega)
					\end{gather*}
					
					is open.
				\end{theorem}
				\begin{proof}
					内容...
				\end{proof}
		\section{Singularities}
			\subsection{Classification}
				\subsubsection{Isolated Singularities}
				\begin{Def}
					A function $f$ is has an \textbf{isolated singularity} at $a\in \CC$, if there exists $R>0$ with
					\begin{itemize}
						\item $f$ is analytic on $B(a,R)-a$;
						\item $f$ is not analyti on $B(a,r)$;
					\end{itemize}
					
					Classification of Isolated Singularities:
					\begin{itemize}
						\item $f$ is called \textbf{removable} if 
						\item $f$ is called \textbf{text} if 
					\end{itemize}
				\end{Def}
				
			\subsection{Residue's Theorem}
				\begin{theorem}
					(Residue's Theorem)
				\end{theorem}
				\begin{proof}
					内容...
				\end{proof}
		\section{The Maximum and Minimum Modulus Theorem}
			\subsection{Maximum Modulus Theorem}
			\begin{theorem}
				(Maximum Modulus Theorem,the First Version) Suppose $G$ be a region, with $f$ analytic on $G$. $a\in G$ with 
				\begin{gather*}
					|f(z)|\leq |f(a)|\qquad z\in G
				\end{gather*}
				
				Then, $f$ is constant.
			\end{theorem}
			\begin{proof}
				内容...
			\end{proof}
			
			\begin{theorem}
				(Maximum Modulus Theorem, the Second Version)
			\end{theorem}
			\begin{proof}
				内容...
			\end{proof}
			
			\begin{theorem}
				(Maximum Modulus Theorem, the Third Version)
			\end{theorem}
			\begin{proof}
				内容...
			\end{proof}
			
			Another proof can also be given as following:
			
			\textbf{Statement} The harmonic function has the maximum/minimum on the boundary.
			
			
			\subsection{Minimum Modulus Theorem}
			\begin{theorem}
				内容...
			\end{theorem}
			\subsection{Schwarz's Theorem}
			\subsection{Hadamard's Three Circle Theorem}
			
		\section{Compactness and Convergence in the Space of Analytic Functions}
			\subsection{The Riemann's Mapping Theorem}
				\begin{theorem}
					内容...
				\end{theorem}
				\begin{proof}
					内容...
				\end{proof}
				
				More topics will be given in: Theory of Univalent Functions.
				\subsubsection{Topic:Biberbach's Polynomial}
				\begin{lemma}
					内容...
				\end{lemma}
				\begin{proof}
					内容...
				\end{proof} 
			\section{Weierstrass's Factorization Theorem}
				\begin{theorem}
					内容...
				\end{theorem}
				\begin{proof}
					内容...
				\end{proof}
		\section{Runge's Theorem}
		\section{Analytic Continuation}
			Original Riemann Surface: Solution of algebraic function
			\begin{gather*}
				f(w,z)=0
			\end{gather*}
			
			Especially: If $f(w_0,z_0)=0;f_z(w_0,z_0)\not=0\Rightarrow $ Locally $z=z(w)$.
			
			\begin{Def}
				\begin{itemize}
					\item Suppose $\gamma:[0,1]\rightarrow \Omega$, then $D_j(\gamma(t_i),r_i)\subseteq \Omega$ is a \textbf{chain of disks} iff
					\begin{itemize}
						\item There exists cut $0=t_1<\dots<t_N=1$,
						\begin{gather*}
							\gamma([t_i,t_{i+1}])\subseteq D_i\cap D_{i+1}
						\end{gather*}
					\end{itemize}
					\item $g\in \mathcal{H}(V)$ is an \textbf{analytic continuation} of $f$ if there exists a chain of disks with $f_j=f_{j+1}$ on $D_j\cap D_{j+1}$ and 
					\begin{gather*}
						f_0=f\qquad f_J=g
					\end{gather*}
				\end{itemize}
			\end{Def}
			
			\begin{lemma}
				The analytic continuation $g$ of $f$ along $\gamma$ only depends of $f$ and $\gamma$, but not on the choice of circles. 
				
				In particular, it is unique.
			\end{lemma}
			\begin{proof}
				内容...
			\end{proof}
			
			\begin{theorem}
				Suppose $\gamma_0,\gamma_1$ two homotopic curves relative to some region $\Omega\subseteq \CC$, with initiao and end points $p,q$. 
				
				Let $U$ be a neighborhood of $p$, assume $f\in \mathcal{H}(U)$, then the analytic continuation agrees locally around $p,q$.
			\end{theorem}
			\begin{proof}
				内容...
			\end{proof}
			
			\textbf{Remark}
			
			\subsubsection{Example: $\Gamma(z)$}
			
			Consider
			\begin{gather*}
				\Gamma(z):=\int_0^\infty e^{-t}t^{z-1}dt
			\end{gather*}
			
			It is analytic on the half plane: 
	\chapter{Harmonic Functions}
	\chapter{Entire Functions}
		\section{Convergence, and Normal Families}
			\begin{lemma}
				Suppose
				\begin{gather*}
					\{f_n\}_{n=1}^\infty \subseteq \mathcal{H}(\Omega)
				\end{gather*}
				
				converges uniformly on compact subsets of $\Omega$ to a function $f$, then 
				\begin{itemize}
					\item $f\in \mathcal{H}(\Omega)$;
					\item $f_n^{(k)}\rightarrow f^{(k)}$ uniformly on compact subset of $\Omega$ for each $k\geq 1$. 
				\end{itemize}
				
				Furthermore, suppose 
				\begin{gather*}
					\sup_{n\geq 1}\#\{z\in \Omega|f_n(z)=w\}\leq N\leq \infty
				\end{gather*}
				
				Then, either $f\equiv w$ in $\Omega$ or $\#\{z\in \Omega|f(z)=w\}\leq N$.
				
				\textbf{Remark} The cardinalities include multiplicity.
			\end{lemma}
			\begin{proof}
				内容...
			\end{proof}
			
			\begin{Def}
				(Normal Family) One says that $\mathcal{F}=\{f_\alpha\}_{\alpha\in A}\subseteq \mathcal{H}(\Omega)$ is a \textbf{normal family} provided for each compact $K\subseteq \Omega$ one has
				\begin{gather*}
					\sup_{z\in K}\sup_{\alpha\in \mathcal{A}}|f_\alpha(z)|<\infty
				\end{gather*} 
			\end{Def}
			
			\begin{proposition}
				One says: $\mathcal{F}=\{f_\alpha\}_{\alpha\in \mathcal{A}}\subseteq \mathcal{H}(\Omega)$ is a \textbf{normal family}. Then, there exists a sequence $f_n$ in $\mathcal{F}$ that converges uniformly on compact subset of $\Omega$. The limit again being a holomorphic function.
			\end{proposition}
			\begin{proof}
				内容...
			\end{proof}
			
			\begin{theorem}
				Let $\mathscr{F}$ be a family of meromorphic functions in $\Omega$, such that for each $z_0\in \Omega$ there is a neighborhood $U(z_0)\subseteq \Omega$ and constant $M(z_0)$ such that:
				\begin{itemize}
					\item Either $|f(z)|<M(z_0)$;
					\item Or $\frac{1}{|f(z)|}<M(z_0),\forall z\in U(z_0),\forall f\in \mathscr{F}$
				\end{itemize}
				
				Then, $\mathscr{F}$ is normal.
			\end{theorem}
			\begin{proof}
				内容...
			\end{proof}
		\section{Mittag-Lefflers Theorem}
			Review: Mittag-Leffers system,HA
			
			\begin{theorem}
				(Mittag-Lefflers Theorem) Given $\{z_n\}_{n=1}^N\subseteq \CC$, with $|z_n|\rightarrow \infty$ if $N\rightarrow \infty$, and polynomial $P_n$ with positive degrees and no constant terms, there exists $f\in \mathcal{M}(\CC)$ so that $f$ has poles exactly at $z_n$, and
				\begin{gather*}
					f(z)-P_n(\frac{1}{z-z_n})
				\end{gather*} 
				
				is analytic around $z_n$ for each $1\leq n\leq N$.
			\end{theorem}
			\begin{proof}
				内容...
			\end{proof}
			
			\begin{corollary}
				Consider the following exact sequence
				\begin{gather*}
					0\rightarrow \mathcal{O}\rightarrow \mathcal{M}\rightarrow \mathcal{M}/\mathcal{O}\rightarrow 0
				\end{gather*}
				
				which induces a long sequence:$$0 \to H^0(X, \mathcal{O}) \xrightarrow{i} H^0(X, \mathcal{M}) \xrightarrow{\pi} H^0(X, \mathcal{M}/\mathcal{O}) \xrightarrow{\delta} H^1(X, \mathcal{O}) \to \dots$$
				
				where $X\subseteq\CC$ is a domain.Now, we have:$$H^1(X, \mathcal{O}) = 0$$
			\end{corollary}
			\begin{proof}
				内容...
			\end{proof}
		\section{Runge's Theorem}
			\subsection{Runge's Theorem}
			Firstly: If $f\in \mathcal{O}(U)$, then it is analytic, and has an approxiamation
			\begin{gather*}
				f=\lim_{n}\sum a_i (z-z_0)^i;z\in D(z_0,r)
			\end{gather*}
			
			But a question is:
			\begin{center}
				If we can replace disk $D(z_0,r)$ with \textbf{arbitrary region}?
			\end{center}
			
			The answer is \textbf{no}: 
			\begin{example}
				Consider $G=\{z\in \CC|0<|z|<1\}$. Then, if $p_n\rightarrow f$, then consider the path
				\begin{gather*}
					\gamma=e^{it}
				\end{gather*}
				
				We have 
				\begin{gather*}
					\lim \oint_\gamma p_n=\oint_\gamma \lim p_n=\oint_\gamma f=0
				\end{gather*}
				
				But, $\oint z^{-1}=0$, and $z^{-1}$ is analytic on $G$. The main reason is simply-coonectedness.
			\end{example}
			
			\begin{lemma}
				Suppose $K$ a \underline{compact} subset of the region $G$, then there are upper line segments
				\begin{gather*}
					\gamma_1,\dots,\gamma_n\subseteq G-K
				\end{gather*}
				
				such that for every function $f$ in $G$,
				\begin{gather*}
					f(z)=\sum \frac{1}{2\pi i}\int_{\gamma_k} \frac{f(w)}{w-z}dw
				\end{gather*}
			\end{lemma}
			\begin{proof}
				内容...
			\end{proof}
			
			\begin{lemma}
				$B(E)$ is a closed subalgebra of $C(K,\CC)$ that contains every rational function with a pole in $E$.
			\end{lemma}
			\begin{proof}
				内容...
			\end{proof}
			
			\begin{lemma}
				内容...
			\end{lemma}
			\begin{proof}
				内容...
			\end{proof}
			
			\begin{lemma}
				内容...
			\end{lemma}
			\begin{proof}
				内容...
			\end{proof}
			
			\begin{theorem}
				Let $K$ compact subset of $\CC$, and $E$ a subset of $\CC_\infty-K$. 
				
				If: 
				\begin{itemize}
					\item $E$ meets each (connected) components of $\CC_\infty-K$;
					\item $f$ is analytic in an open set containing $K$;
					\item $\epsilon>0$
				\end{itemize}
				
				Then, there exists a rational number $R(z)$, whose only pole lie in $E$, and such that
				\begin{gather*}
					|f(z)-R(z)|<\epsilon
				\end{gather*}
			\end{theorem}
			\begin{proof}
				内容...
			\end{proof}
			
			
			\textbf{Remark} Here, we could give another proofs:
			\begin{itemize}
				\item 
			\end{itemize}
			\subsection{The Density of Functions}
			\subsection{Polynomially Convex}
			\begin{Def}
				内容...
			\end{Def}
			
			
	\chapter{The Range of Analytic Functions}
	\chapter{Univalent Functions}
		Given in a new topic of part.
	
	\newpage
	
	Continuation of Analytic Functions, given in Riemann Surface.
\end{document}